% Erzeugt mit Trendo OS – mit LuaLaTeX oder XeLaTeX übersetzen
\documentclass[11pt,a4paper]{article}
\usepackage{fontspec}
\usepackage{amsmath,amssymb}
\usepackage[ngerman]{babel}
\usepackage[margin=2cm]{geometry}
\title{aufgaben-libreoffice}
\date{}
\begin{document}
\maketitle
\section*{Löse nach x: 5x − 3 = 2x + 9}
\textit{Gleichung lösen}

\paragraph{Lösungsweg}
\begin{enumerate}
  \item Alles auf eine Seite bringen, ausmultiplizieren und zusammenfassen\\
  \texttt{5x − 3 = 2x + 9  ⇔  3x − 12 = 0}
  \item +12 auf beiden Seiten\\
  \texttt{3x = 12}
  \item Beide Seiten durch 3 teilen\\
  \texttt{x = 12 : 3 = 4}
\end{enumerate}
\paragraph{Ergebnis}
\[ x_{1} = 4 \]
\paragraph{Probe}
\begin{itemize}
  \item ✓ x = 4 einsetzen: links 17, rechts 17 – stimmt (exakt).
  \item ✓ Unabhängige Zählung mit der Sturm-Kette: genau 1 verschiedene reelle Nullstelle(n) – alle gefunden.
\end{itemize}
\section*{Löse nach x: x² + 2x − 15 = 0}
\textit{Gleichung lösen}

\paragraph{Lösungsweg}
\begin{enumerate}
  \item p-q-Formel mit p = 2 und q = −15: x = −p/2 ± √((p/2)² − q)\\
  \texttt{x = −1 ± √((1)² − (−15)) = −1 ± √(16)}
  \item √(16) = 4 – zwei Lösungen\\
  \texttt{x₁ = −1 + 4 = 3;  x₂ = −1 − 4 = −5}
\end{enumerate}
\paragraph{Ergebnis}
\[ x_{1} = -5,\quad x_{2} = 3 \]
\paragraph{Probe}
\begin{itemize}
  \item ✓ x = −5 einsetzen: links 0, rechts 0 – stimmt (exakt).
  \item ✓ x = 3 einsetzen: links 0, rechts 0 – stimmt (exakt).
  \item ✓ Unabhängige Zählung mit der Sturm-Kette: genau 2 verschiedene reelle Nullstelle(n) – alle gefunden.
\end{itemize}
\section*{Löse das Gleichungssystem}
\textit{Gleichungssystem}

\paragraph{Lösungsweg}
\begin{enumerate}
  \item Jede Gleichung ordnen: Variablen links, Zahl rechts\\
  \texttt{I: x + y = 10 | II: x − y = 2}
  \item x aus den anderen Gleichungen entfernen (Additionsverfahren): II − I\\
  \texttt{I: x + y = 10 | II: −2y = −8}
  \item Gleichung II durch −2 teilen, damit vor y eine 1 steht\\
  \texttt{I: x + y = 10 | II: y = 4}
  \item y aus den anderen Gleichungen entfernen (Additionsverfahren): I − II\\
  \texttt{I: x = 6 | II: y = 4}
  \item Lösung ablesen\\
  \texttt{x = 6;  y = 4}
\end{enumerate}
\paragraph{Ergebnis}
\[ x = 6,\quad y = 4 \]
\paragraph{Probe}
\begin{itemize}
  \item ✓ Gleichung I (x + y = 10) mit der Lösung: links 10, rechts 10.
  \item ✓ Gleichung II (x − y = 2) mit der Lösung: links 2, rechts 2.
\end{itemize}
\section*{Berechne: 2/3 + 3/4}
\textit{Bruchrechnen}

\paragraph{Lösungsweg}
\begin{enumerate}
  \item Hauptnenner finden: kgV(3, 4) = 12; beide Brüche darauf erweitern\\
  \texttt{2/3 + 3/4 = 8/12 + 9/12}
  \item Zähler addieren\\
  \texttt{= 17/12}
\end{enumerate}
\paragraph{Ergebnis}
\[ \frac{17}{12} \]
\paragraph{Probe}
\begin{itemize}
  \item ✓ Gegenrechnung mit dem Taschenrechner (Gleitkomma): 1,4166666667 – Bruch 17/12 ergibt 1,4166666667.
\end{itemize}
\section*{Faktorisiere: x² − 16}
\textit{Faktorisieren}

\paragraph{Lösungsweg}
\begin{enumerate}
  \item 3. binomische Formel rückwärts: a² − b² = (a + b)(a − b)\\
  \texttt{x² − 16 = (x + 4)(x − 4)}
  \item Ergebnis\\
  \texttt{x² − 16 = (x + 4)(x − 4)}
\end{enumerate}
\paragraph{Ergebnis}
\[ \left(x + 4\right) \cdot \left(x - 4\right) \]
\paragraph{Probe}
\begin{itemize}
  \item ✓ Ausmultipliziert ergibt (x + 4)(x − 4) wieder x² − 16 – genau der Ausgangsterm.
\end{itemize}
\section*{Integriere: ∫ (4x³ − 2x) dx}
\textit{Integrieren}

\paragraph{Lösungsweg}
\begin{enumerate}
  \item Summenregel: jedes Glied einzeln integrieren\\
  \texttt{∫ (4x³ − 2x) dx = ∫ 4x³ dx − ∫ 2x dx}
  \item Faktorregel: der Faktor 4 kommt vor das Integral\\
  \texttt{∫ 4x³ dx = 4 · ∫ x³ dx}
  \item Potenzregel: ∫ uⁿ du = uⁿ⁺¹/(n + 1) mit n = 3\\
  \texttt{∫ x³ dx = x⁴/4}
  \item Faktorregel: der Faktor −2 kommt vor das Integral\\
  \texttt{∫ −2x dx = −2 · ∫ x dx}
  \item Potenzregel: ∫ xⁿ dx = xⁿ⁺¹/(n + 1) mit n = 1\\
  \texttt{∫ x dx = x²/2}
  \item Stammfunktion (C ist eine beliebige Konstante)\\
  \texttt{∫ (4x³ − 2x) dx = x⁴ − x² + C}
\end{enumerate}
\paragraph{Ergebnis}
\[ \int 4 \cdot x^{3} - 2 \cdot x\,dx = x^{4} - x^{2} + C \]
\paragraph{Probe}
\begin{itemize}
  \item ✓ Probe durch Ableiten: F' stimmt an 6 Stellen mit f überein – bei x = 0,7: F' ≈ −0,028, f = −0,028.
\end{itemize}
\section*{Leite ab: f(x) = x²·e\textasciicircum{}x}
\textit{Ableiten}

\paragraph{Lösungsweg}
\begin{enumerate}
  \item Produktregel: (u·v)' = u'·v + u·v' mit u = x², v = e\textasciicircum{}x\\
  \texttt{(x²·e\textasciicircum{}x)' = x²·e\textasciicircum{}x + 2x·e\textasciicircum{}x}
  \item Potenzregel: (xⁿ)' = n·xⁿ⁻¹ mit n = 2\\
  \texttt{(x²)' = 2x}
  \item e-Funktion: (eᵘ)' = eᵘ·u'\\
  \texttt{(e\textasciicircum{}x)' = e\textasciicircum{}x}
  \item Ableitung von x\\
  \texttt{(x)' = 1}
  \item Zusammenfassen\\
  \texttt{f'(x) = x²·e\textasciicircum{}x + 2x·e\textasciicircum{}x}
\end{enumerate}
\paragraph{Ergebnis}
\[ f'(x) = x^{2} \cdot e^{x} + 2 \cdot x \cdot e^{x} \]
\paragraph{Probe}
\begin{itemize}
  \item ✓ Gegenrechnung mit dem Differenzenquotienten an 6 Stellen – bei x = 0,7: Differenzenquotient 3,805993, Ableitung 3,805993.
\end{itemize}
\end{document}
